% !TEX root = ../main.tex

\section{Our Rust Implementation}
\label{app:zq-rust}

\textbf{Software.} We use arkworks 0.4 in a Cargo workspace, with \texttt{sha2} for hashing and optional \texttt{rayon} support for parallel execution. Table 13 gives the size and responsibilities of the crates.

\textbf{Table 13.} Crates of our Rust implementation. Lines exclude blanks and comments.

\begingroup
\setlength{\tabcolsep}{2pt}\small
\begin{longtable}{@{}>{\raggedright\arraybackslash}p{0.2520\linewidth}>{\raggedright\arraybackslash}p{0.5040\linewidth}>{\raggedright\arraybackslash}p{0.1440\linewidth}@{}}
\toprule
Crate & Content & Lines \\
\midrule\endfirsthead
\toprule
Crate & Content & Lines \\
\midrule\endhead
\bottomrule\endfoot
\texttt{zeequent-types} & Parameters, orders, the plaintext book, the receipt & 237 \\
\texttt{zeequent-kzg} & Commitments, openings, transcript, setup import, file codec & 754 \\
\texttt{zeequent-core} & Clearing rule, constraints, prover, verifier, committed book, update validity & 2,596 \\
\texttt{zeequent-cli} & The command-line tool & 1,090 \\
Tests & 44 tests & 1,459 \\
\end{longtable}
\endgroup

\textbf{Setup.} The benchmarks use locally generated test parameters. For an external setup, we read the reference string from a powers-of-tau ceremony file: \([\tau^i]_1\) for \(i \le 2M + 8\), \([1]_2\), \([\tau]_2\), and the file's \([\alpha\tau^i]_1\) section as the blinding basis, so \(\gamma\) is the ceremony's \(\alpha\). From the powers we derive what depends on the domain: the Lagrange-basis points \([L_i(\tau)]_1\), by an inverse FFT over group elements, and \([Z_H(\tau)]_2\). A verifier that receives a setup from another party checks with pairings that its points are powers of one \(\tau\) and that the Lagrange points match them. The check does not show that \(\tau\) is unknown; that rests on the ceremony.

\textbf{Commitments.} We commit directly to the evaluations in each witness column. In the Lagrange basis, \([p(\tau)]_1=\sum_i p[i]\,[L_i(\tau)]_1\) requires one size-\(n\) multi-scalar multiplication, with the column entries as scalars. Our earlier coefficient-based implementation took 3.6 times as long to prove. We commit separately to the mask \(\rho(X)Z_H(X)\) and the blinding term \(\gamma r(X)\), then add both commitments. For columns sampled within the proof, the mask and blinding degrees equal the number of opening points. Section 4.1 gives the degree-three masks for the running book commitments.

\textbf{Quotient.} The combined constraint polynomial has degree up to \(3n + 3\), so we evaluate it on a coset of size \(4n\) that avoids \(H\), divide pointwise by \(Z_H\), and interpolate. The large coset is a union of four cosets of \(H\), and we process them one at a time with transforms of size \(n\). Before each transform we reduce a masked column, whose degree slightly exceeds \(n\), modulo \(X^n - s\), where \(s\) is the constant value of \(X^n\) on that coset. On each small coset \(Z_H\) is a constant and \(\omega x\) is the next point, so a shifted column is a rotation. If the constraints do not hold, the quotient has non-zero coefficients above degree \(2n + 8\) and the prover stops.

\textbf{Openings.} To open several polynomials at a point \(z\), the prover combines them with powers of a challenge, divides by \(X - z\), and commits to the quotient \(W\); it sends one combined blinding evaluation \(b\). With \(C\) and \(v\) the combined commitment and value, the verifier checks both opening points with two pairings:

\[
e\Big(\sum_{k} r^k \big(C_k - v_k [1]_1 - b_k\, \Gamma + z_k W_k\big),\ [1]_2\Big) = e\Big(\sum_{k} r^k W_k,\ [\tau]_2\Big),
\]

where \(r\) is a further challenge. We evaluate the left pairing argument with a single multi-scalar multiplication over the proof commitments. This reduced verification time by about a factor of eight relative to combining the commitments pairwise.

\textbf{Public polynomials.} We store a public polynomial as its non-zero values on \(H\). The verifier evaluates it at \(\zeta\) from \(L_i(\zeta) = \omega^i(\zeta^n - 1)/(n(\zeta - \omega^i))\) with one batch inversion, at a cost linear in the number of non-zero values: the width of the plateau for \(Mask_P\), 256 for the range table.

\textbf{Transcript.} We maintain a SHA-256 digest. To absorb a message, we hash the current digest with the message label, length, and contents. To derive a challenge, we reduce 64 derived bytes modulo the field size and hash the challenge back into the state. The protocol label, parameters, receipt, and ordered commitments bind each challenge to the public statement and prior messages.

\textbf{Proof format.} We serialize the column and quotient commitments, followed by two openings. Each opening includes its point, evaluations, combined blinding evaluation, and witness. A compressed curve point occupies 32 bytes. If \(c\) columns are opened and \(s\) also open at \(\omega\zeta\), the proof occupies \(32(2c+s+8)+12\) bytes. With \(c=88\) and \(s=7\) this gives 6,124 bytes. Decompression requires a base-field square root per point and accounts for most single-threaded verification time.

\textbf{Determinism.} At the close we derive the prover's randomness from a hash of its private state, which includes the mask and blinding of every book column. Closing the same batch twice produces the same bytes, and the tool refuses orders after a close. This implements the one-proof condition of Appendix~\ref{app:zq-secrecy-proof} for an honest exchange.

\textbf{Batch files.} The command-line tool keeps a batch in a directory (Table 14) and has one command per protocol step: \texttt{setup}, \texttt{open}, \texttt{order}, \texttt{check-receipt}, \texttt{close}, \texttt{verify}, and \texttt{fill}.

\textbf{Table 14.} Files of one batch, for 32-bit volumes.

\begingroup
\setlength{\tabcolsep}{2pt}\small
\begin{longtable}{@{}>{\raggedright\arraybackslash}p{0.1800\linewidth}>{\raggedright\arraybackslash}p{0.3690\linewidth}>{\raggedright\arraybackslash}p{0.1530\linewidth}>{\raggedright\arraybackslash}p{0.1980\linewidth}@{}}
\toprule
File & Content & Public & Size \\
\midrule\endfirsthead
\toprule
File & Content & Public & Size \\
\midrule\endhead
\bottomrule\endfoot
Opening & Four commitments with their emptiness proofs & Yes & 1,576 B \\
Board & The log & Yes & 416 B per order \\
Update receipt & Column, tick, quantity, blinding, commitment before and after & To its trader & 192 B \\
Outcome & The clearing receipt, or no-cross & Yes & About 100 B \\
Clearing proof &  & Yes & 6,124 B \\
Update-validity proof &  & Yes & 1,700 B \\
State & Plaintext book, masks, blindings, update history & No & Grows with the orders \\
\end{longtable}
\endgroup

\textbf{Tests.} Table 15 pairs implementation checks with the claims they exercise. In tests with a false statement we disable the prover's local checks, pass its output to the verifier, and check that the verifier rejects it.

\textbf{Table 15.} Claims and the tests that support them.

\begingroup
\setlength{\tabcolsep}{2pt}\small
\begin{longtable}{@{}>{\raggedright\arraybackslash}p{0.3150\linewidth}>{\raggedright\arraybackslash}p{0.5850\linewidth}@{}}
\toprule
Claim & Test \\
\midrule\endfirsthead
\toprule
Claim & Test \\
\midrule\endhead
\bottomrule\endfoot
The code's clearing rule is Definition 2 & 2,000 random books; the receipt is recomputed from the definition by direct search \\
The constraints admit one receipt (Theorem 1) & 300 random books of 6 ticks; every receipt with values up to the total volume is tried, and exactly one is accepted \\
The polynomial system equals the integer constraints & Receipt values and witness cells are changed one at a time and proved with the prover's checks off; the verifier rejects each \\
Every tie-break case proves and verifies & One book per case of the clearing-rule appendix, and a padded domain \\
A crossing book cannot be declared no-cross & Direct test \\
A logged step cannot cancel an order (no removal) & The attack of Section 4.3 is run; no update-validity proof verifies, whatever quantity and tick the exchange claims \\
A repeated close reveals nothing new & The two proofs are compared byte for byte \\
Range checks are exact & An 8-bit quantity in 3-bit limbs, a width that is not a multiple of the limb: values in range pass and values out of range fail \\
\end{longtable}
\endgroup

\section{Our Noir Implementation}
\label{app:zq-noir}

We wrote the Noir circuit directly from the clearing constraints. It supplies an independent implementation of the rule and the general-purpose benchmark in Section~\ref{sec:zq-comparison}.

\textbf{Statement.} The witness consists of four arrays of \(N\) unsigned 32-bit volumes and a salt. The crossing circuit publishes nine field elements (288 bytes) for the book digest and eight receipt values. The no-cross circuit publishes only the digest.

\textbf{Circuit model.} Noir constrains the result of ordinary arithmetic, including the sums used to compute depth. Its unsigned integer types enforce bounds, so a subtraction that would underflow cannot satisfy the circuit. An unconstrained function can supply a witness value, but the circuit must still check that value. The backend implements the polynomial domain, quotient, openings, and lookups.

\textbf{Executable volume.} An unconstrained function proposes the minimum of the two depths. The circuit subtracts this value from each depth in unsigned 64-bit arithmetic, enforcing the differences and their non-negativity (F9, F10, F12, F13). It then checks that their product is zero (F11). Together the checks force the proposed value to equal the smaller depth.

\textbf{Types.} Table 16 maps Noir types to the constraints they enforce. Arrays have exactly \(N\) entries. Native comparisons implement the inequalities in F14 and F19 without an explicit slack column or limb witnesses at the circuit-language level.

\textbf{Table 16.} Types of the circuit and the constraints they replace.

\begingroup
\setlength{\tabcolsep}{2pt}\small
\begin{longtable}{@{}>{\raggedright\arraybackslash}p{0.2070\linewidth}>{\raggedright\arraybackslash}p{0.2250\linewidth}>{\raggedright\arraybackslash}p{0.4680\linewidth}@{}}
\toprule
Quantity & Type & Replaces \\
\midrule\endfirsthead
\toprule
Quantity & Type & Replaces \\
\midrule\endhead
\bottomrule\endfoot
Volume at one tick & 32-bit unsigned & F1, F2, F24, F25 \\
Depths, \(Min\), surpluses, \(\Delta\) & 64-bit unsigned & F12, F13, and the range of every accumulated value \\
\(c\), \(d\), \(p^*\) & 32-bit unsigned &  \\
Side & Boolean &  \\
Array of length \(N\) &  & F3, F4, F26, F27 \\
\end{longtable}
\endgroup

\textbf{Receipt.} An unconstrained function computes a candidate receipt from the book. The circuit then applies the scalar and per-tick checks in Table 17. By Theorem 1, every accepted candidate equals the rule's receipt. The backend handles reads at \(c\), \(d\), \(p^*\), and their neighbours as memory lookups because those positions depend on the witness.

\textbf{Table 17.} Steps of the clearing check in our Noir implementation.

\begingroup
\setlength{\tabcolsep}{2pt}\small
\begin{longtable}{@{}>{\raggedright\arraybackslash}p{0.1710\linewidth}>{\raggedright\arraybackslash}p{0.5760\linewidth}>{\raggedright\arraybackslash}p{0.1530\linewidth}@{}}
\toprule
Step & Assertion & Constraints \\
\midrule\endfirsthead
\toprule
Step & Assertion & Constraints \\
\midrule\endhead
\bottomrule\endfoot
Receipt & \(c \le p^* \le d < N\); \(V_{max} \ge 1\); \(V_{min\Delta} \le P\); \(s\), \(P\), \(J\) zero when balanced & Scalar checks \\
Depth & Running sums from the top tick and from the bottom tick & F5--F8, F28--F31 \\
Each tick & The minimum; \(\Delta = SurpA + SurpB\); \(Min \le V_{max}\); \(\Delta \ge V_{min\Delta}\) on \([c,d]\) & F9--F14, F18, F19 \\
Plateau & \(Min[c] = Min[d] = V_{max}\); \(Min\) strictly lower at \(c - 1\) and \(d + 1\) where they exist & F15--F17 \\
Clearing tick & The surpluses at \(p^*\) are \(V_{min\Delta}\) on the long side and 0 on the other & F20 \\
Tie-break & If \(p^* > c\) and supply is long or the batch is balanced, \(\Delta[p^* - 1] > V_{min\Delta}\) & F21 \\
Rationing & \(P\) is the long side's volume at \(p^*\); \(J\) is the short side's eligible IO depth at \(p^*\) & F32, F33 \\
\end{longtable}
\endgroup

\textbf{Digest.} We pack the four volumes at each tick into \(b+2^{32}a+2^{64}b^{io}+2^{96}a^{io}\). A width-four Poseidon2 sponge with rate three and capacity one absorbs the resulting field elements. We initialize the state to \((\text{salt},0,0,N)\), permute once before absorption, and take the first state element as the output after absorption. The private salt prevents an observer from testing a guessed low-volume book against the digest without also guessing the salt. At 1,000 ticks a Pedersen hash of the unpacked volumes required 1.78 million gates. Our sponge uses about 25,000. This arrangement of the standard permutation has not received a separate security analysis.

\textbf{Size.} The circuit has 12,757 gates at 100 ticks, 100,357 at 1,000, and 197,666 at 2,000, about 100 per tick. The proof is 14,656 bytes and the verification key 3,680 bytes at every size. The number of ticks is a compile-time constant, so each book size has its own verification key.

\textbf{Tests.} Four of 20 library tests check hand-computed receipts for a demand-long book, a supply-long book, a balanced book, and a book with one maximizing tick and IO interest. Twelve tests submit false receipts and check which constraint rejects them. The remaining tests cover no-cross and the dependence of the digest on the book and salt. Across the 57 calibrated and replayed books and the 362 Cboe-derived books, the six fields exported by our Rust benchmark agree with our Noir implementation. Independent recomputation from the input volumes checks all eight fields returned by our Noir implementation, including the two rationing quantities.

\textbf{Toolchain.} We compile with \texttt{nargo} 1.0.0 and prove with Barretenberg 6.0.0 in its zero-knowledge UltraHonk mode with a Poseidon2 transcript. The circuit is 208 lines and its tests 161.
